\chapter*{Abstract}

In this thesis simulations of \acrlong{tbl} across curved walls under \acrlong{zpg} conditions were performed to enable an isolated study of curvature induced changes in the \acrshort{tbl}.
In a curved geometry a streamwise pressure gradient occurs whose effects superimpose the curvature induced effects on the \acrlong{tbl}.
Therefore, the pressure field is manipulated to generate \acrshort{zpg} conditions in the entire simulated domain.
In this project a new approach to build a curved setup with \acrshort{zpg} conditions was tested.
It is based on an optimisation algorithm presented by \citet{morita_applying_2022}, which performs a shape optimisation  of the upper wall of the simulation domain in order to match a prescribed distribution of the Clauser parameter $\beta$ at the lower wall.
The resulting velocity field is extracted and imposed as a \benennung{Dirichlet} condition on the upper wall of a \acrlong{dns} setup.
With this setup \acrlong{dns} of \acrlong{tbl} in three convex curved domains with different geometries were performed.
The results were compared with curved simulations under reference conditions and reference data from flat-plate \acrlong{tbl} simulations.
The obtained results obtained demonstrate the capability of the optimisation algorithm to manipulate the pressure field towards a constant distribution on a global level, but shows limitations when it comes to extreme curvature situations as a 90° bend.
